\subsection{Filters, Ultrafilters, and Sets of Reals}

\begin{definition} \label{def:filters-ideals-and-duality} \hlt{Filters, Ideals, and Duality}\\
Let $X$ be a set, let $\mathcal{F},\mathcal{I}\subseteq\mathcal{P}(X)$, and let $A\subseteq X$.
\begin{enumerate}[label=(\roman*)]
\setlength{\itemsep}{0pt}
\item The family $\mathcal{F}$ is a \hlt{filter} on $X$ iff $X\in\mathcal{F}$, $\emptyset\notin\mathcal{F}$, $(\forall B,C\in\mathcal{F})\,B\cap C\in\mathcal{F}$, and $(\forall B\in\mathcal{F})(\forall C\subseteq X)\,[B\subseteq C\Rightarrow C\in\mathcal{F}]$.
\item The family $\mathcal{I}$ is an \hlt{ideal} on $X$ iff $\emptyset\in\mathcal{I}$, $X\notin\mathcal{I}$, $(\forall B,C\in\mathcal{I})\,B\cup C\in\mathcal{I}$, and $(\forall B\in\mathcal{I})(\forall C\subseteq B)\,C\in\mathcal{I}$.
\item If $\mathcal{F}$ is a filter, its \hlt{dual ideal} is $\mathcal{I}_{\mathcal{F}}\defeq\{X\setminus B:B\in\mathcal{F}\}$. If $\mathcal{I}$ is an ideal, its \hlt{dual filter} is $\mathcal{F}_{\mathcal{I}}\defeq\{X\setminus B:B\in\mathcal{I}\}$.
\item If $\mathcal{I}$ is an ideal, then $A$ is \hlt{$\mathcal{I}$-positive} iff $A\notin\mathcal{I}$.
\end{enumerate}
\end{definition}

\begin{definition} \label{def:generated-principal-and-frechet-filters} \hlt{Generated, Principal, and Fr\'echet Filters}\\
Let $X$ be a set and let $\mathcal{B}\subseteq\mathcal{P}(X)$ satisfy $\mathcal{B}\neq\emptyset$, $\emptyset\notin\mathcal{B}$, and $(\forall B_0,B_1\in\mathcal{B})\,B_0\cap B_1\in\mathcal{B}$.
\begin{enumerate}[label=(\roman*)]
\setlength{\itemsep}{0pt}
\item The \hlt{filter generated by $\mathcal{B}$} is $\langle\mathcal{B}\rangle\defeq\{A\subseteq X:(\exists B\in\mathcal{B})\,B\subseteq A\}$.
\item For $\emptyset\neq D\subseteq X$, the \hlt{principal filter generated by $D$} is $\mathcal{F}_D\defeq\{A\subseteq X:D\subseteq A\}$.
\item For infinite $X$, the \hlt{Fr\'echet filter}, or \hlt{cofinite filter}, is $\operatorname{Fr}(X)\defeq\{A\subseteq X:X\setminus A\text{ is finite}\}$.
\end{enumerate}
\end{definition}

\begin{definition} \label{def:kappa-complete-filter} \hlt{$\kappa$-Complete Filter}\\
Let $X$ be a set, let $\kappa$ be an infinite cardinal, and let $\mathcal{F}$ be a filter on $X$.\\
The filter $\mathcal{F}$ is \hlt{$\kappa$-complete} iff $(\forall\lambda<\kappa)(\forall A\in\mathcal{F}^{\lambda})\,\bigcap\limits_{\xi<\lambda}A(\xi)\in\mathcal{F}$.
\end{definition}

\begin{proposition} \label{prop:basic-filter-and-ideal-constructions} \hlt{Basic Filter and Ideal Constructions}\\
Let $X$ be a set, let $\mathcal{F}$ be a filter on $X$, let $\mathcal{I}$ be an ideal on $X$, and let $\emptyset\neq D\subseteq X$.\\
Let $\mathcal{B}\subseteq\mathcal{P}(X)$ satisfy $\mathcal{B}\neq\emptyset$, $\emptyset\notin\mathcal{B}$, and $(\forall B_0,B_1\in\mathcal{B})\,B_0\cap B_1\in\mathcal{B}$.
\begin{enumerate}[label=(\roman*)]
\setlength{\itemsep}{0pt}
\item $\mathcal{I}_{\mathcal{F}}$ is an ideal and $\mathcal{F}_{\mathcal{I}}$ is a filter on $X$.
\item $\langle\mathcal{B}\rangle$ is the least filter containing $\mathcal{B}$.
\item $\mathcal{F}_D$ is a filter; if $X$ is infinite, then $\operatorname{Fr}(X)$ is a filter whose dual ideal is the family of finite subsets of $X$.
\item The filter $\mathcal{F}$ is $\aleph_0$-complete.
\end{enumerate}
\begin{proof}
\begin{enumerate}[label=(\roman*)]
\setlength{\itemsep}{0pt}
\item The equalities $X\setminus X=\emptyset$ and $X\setminus\emptyset=X$ show that $\mathcal{I}_{\mathcal{F}}$ contains $\emptyset$ but not $X$. If $A,B\in\mathcal{F}$, then $(X\setminus A)\cup(X\setminus B)=X\setminus(A\cap B)\in\mathcal{I}_{\mathcal{F}}$.\\
Let $A\in\mathcal{F}$ and $E\subseteq X\setminus A$. Then $A\subseteq X\setminus E$, so $X\setminus E\in\mathcal{F}$ and $E\in\mathcal{I}_{\mathcal{F}}$. Thus $\mathcal{I}_{\mathcal{F}}$ is an ideal.\\
The same two complement equalities give $X\in\mathcal{F}_{\mathcal{I}}$ and $\emptyset\notin\mathcal{F}_{\mathcal{I}}$. If $A,B\in\mathcal{I}$, then $(X\setminus A)\cap(X\setminus B)=X\setminus(A\cup B)\in\mathcal{F}_{\mathcal{I}}$.\\
Let $A\in\mathcal{I}$ and $X\setminus A\subseteq E\subseteq X$. Then $X\setminus E\subseteq A$, so $X\setminus E\in\mathcal{I}$ and $E\in\mathcal{F}_{\mathcal{I}}$. Hence $\mathcal{F}_{\mathcal{I}}$ is a filter.
\item Choose $B_0\in\mathcal{B}$. Then $B_0\subseteq X$ gives $X\in\langle\mathcal{B}\rangle$. Since no member of $\mathcal{B}$ is contained in $\emptyset$, one has $\emptyset\notin\langle\mathcal{B}\rangle$.\\
If $A,A'\in\langle\mathcal{B}\rangle$, choose $B,B'\in\mathcal{B}$ with $B\subseteq A$ and $B'\subseteq A'$. Then $B\cap B'\in\mathcal{B}$ and $B\cap B'\subseteq A\cap A'$.\\
Thus $\langle\mathcal{B}\rangle$ is closed under intersections, and upward closure follows from transitivity of inclusion. The inclusion $\mathcal{B}\subseteq\langle\mathcal{B}\rangle$ holds.\\
If a filter $\mathcal{G}$ contains $\mathcal{B}$ and $A\in\langle\mathcal{B}\rangle$, choose $B\in\mathcal{B}$ with $B\subseteq A$. Then $A\in\mathcal{G}$, so $\langle\mathcal{B}\rangle\subseteq\mathcal{G}$.
\item Since $\mathcal{F}_D=\langle\{D\}\rangle$, the family $\mathcal{F}_D$ is a filter by (ii). Suppose $X$ is infinite. The set $X$ is cofinite, whereas $\emptyset$ is not. The intersection of two cofinite sets is cofinite because its complement is the union of two finite sets, and every superset of a cofinite set is cofinite. Thus $\operatorname{Fr}(X)$ is a filter. Its complement transform is $\{A\subseteq X:A\text{ is finite}\}$.
\item Let $\lambda<\aleph_0$ and $A\in\mathcal{F}^{\lambda}$. The ordinal $\lambda$ is finite because $\aleph_0$ is the least infinite cardinal. The empty intersection is $X\in\mathcal{F}$, and adjoining one $A(\xi)$ preserves membership by binary closure. Hence the required intersection belongs to $\mathcal{F}$ by induction on $\lambda$.
\end{enumerate}
\hlt{Structural Commentary}: Complementation proves duality; base closure and finite induction handle the filter constructions.\\
\hlt{Key Insight}: Upward closure turns finite-intersection bases into filters.
\end{proof}
\end{proposition}

\begin{definition} \label{def:ultrafilter-and-principality} \hlt{Ultrafilters and Principal Ultrafilters}\\
Let $X$ be a set and let $\mathcal{U}$ be a filter on $X$.
\begin{enumerate}[label=(\roman*)]
\setlength{\itemsep}{0pt}
\item The filter $\mathcal{U}$ is an \hlt{ultrafilter} iff $(\forall A\subseteq X)\,\#(\mathcal{U}\cap\{A,X\setminus A\})=1$.
\item The ultrafilter $\mathcal{U}$ is \hlt{principal} iff $(\exists x\in X)\,\mathcal{U}=\{A\subseteq X:x\in A\}$; otherwise it is \hlt{nonprincipal}.
\end{enumerate}
\end{definition}

\begin{theorem} \label{thm:ultrafilters-are-maximal-filters} \hlt{Ultrafilters Are Exactly the Maximal Filters}\\
Let $X$ be a set and let $\mathcal{F}$ be a filter on $X$. Then $\mathcal{F}$ is an ultrafilter iff $(\forall\mathcal{G}\subseteq\mathcal{P}(X))\,[\mathcal{G}\text{ is a filter on }X\land\mathcal{F}\subseteq\mathcal{G}\Rightarrow\mathcal{G}=\mathcal{F}]$.
\begin{proof}
\begin{enumerate}[label=\roman*.]
\setlength{\itemsep}{0pt}
\item[$(\Rightarrow)$] Suppose $\mathcal{F}$ is an ultrafilter and let $\mathcal{G}$ be a filter with $\mathcal{F}\subseteq\mathcal{G}$\\
If $A\in\mathcal{G}\setminus\mathcal{F}$, then $X\setminus A\in\mathcal{F}\subseteq\mathcal{G}$, so $\emptyset=A\cap(X\setminus A)\in\mathcal{G}$, contradicting properness. Thus $\mathcal{G}=\mathcal{F}$.
\item[$(\Leftarrow)$] Suppose the displayed maximality condition holds, and fix $A\subseteq X$. Assume $A,X\setminus A\notin\mathcal{F}$, and set $\mathcal{B}_A\defeq\{B\cap A:B\in\mathcal{F}\}$.\\
The relation $A=X\cap A$ gives $\mathcal{B}_A\neq\emptyset$, and $(B\cap A)\cap(C\cap A)=(B\cap C)\cap A\in\mathcal{B}_A$ for $B,C\in\mathcal{F}$.\\
If $\emptyset\in\mathcal{B}_A$, choose $B\in\mathcal{F}$ with $B\cap A=\emptyset$. Then $B\subseteq X\setminus A$, so $X\setminus A\in\mathcal{F}$, a contradiction.\\
Hence $\langle\mathcal{B}_A\rangle$ is a filter by Proposition \ref{prop:basic-filter-and-ideal-constructions}(ii). For $B\in\mathcal{F}$, the relation $B\cap A\subseteq B$ gives $B\in\langle\mathcal{B}_A\rangle$, so $\mathcal{F}\subseteq\langle\mathcal{B}_A\rangle$.\\
Since $A\in\mathcal{B}_A\subseteq\langle\mathcal{B}_A\rangle$ but $A\notin\mathcal{F}$, the containment is proper, contradicting maximality.\\
Thus $A\in\mathcal{F}$ or $X\setminus A\in\mathcal{F}$; properness prevents both, so $\mathcal{F}$ is an ultrafilter.
\end{enumerate}
\hlt{Structural Commentary}: One direction excludes proper extensions; the other adjoins any undecided set through a filter base.\\
\hlt{Key Insight}: Maximality is equivalent to deciding every complementary pair.
\end{proof}
\end{theorem}

\begin{principle} \label{princ:ultrafilter-lemma} \hlt{Ultrafilter Lemma}\\
Let $X$ be a set and let $\mathcal{F}$ be a filter on $X$. The Ultrafilter Lemma asserts $(\exists\mathcal{U}\subseteq\mathcal{P}(X))\,[\mathcal{U}\text{ is an ultrafilter on }X\land\mathcal{F}\subseteq\mathcal{U}]$.
\end{principle}

\begin{theorem} \label{thm:ultrafilter-extension-under-zorn} \hlt{Ultrafilter Extension under Zorn's Lemma}\\
Principle \ref{princ:zorns-lemma} implies Principle \ref{princ:ultrafilter-lemma}.
\begin{proof}
\begin{enumerate}[label=(\roman*)]
\setlength{\itemsep}{0pt}
\item \hlt{Extension poset}: Assume Zorn's Lemma, let $X$ be a set, and let $\mathcal{F}$ be a filter on $X$. \\
Let $\mathscr{P}$ be the set of filters on $X$ containing $\mathcal{F}$, ordered by inclusion.\\
Then $\mathcal{F}\in\mathscr{P}$, and inclusion is a partial order on $\mathscr{P}$ by Proposition \ref{prop:inclusion-and-reverse-inclusion-partial-orders}.
\item \hlt{Chain bounds}: The filter $\mathcal{F}$ bounds the empty chain. Let $\mathcal{C}\subseteq\mathscr{P}$ be a nonempty chain and set $\mathcal{G}\defeq\bigcup\limits_{\mathcal{H}\in\mathcal{C}}\mathcal{H}$.\\
Then $X\in\mathcal{G}$, $\emptyset\notin\mathcal{G}$, and $\mathcal{F}\subseteq\mathcal{G}$. If $A,B\in\mathcal{G}$, choose $\mathcal{H}_0,\mathcal{H}_1\in\mathcal{C}$ with $A\in\mathcal{H}_0$ and $B\in\mathcal{H}_1$.\\
Since $\mathcal{C}$ is a chain, one of these filters contains both $A$ and $B$, so $A\cap B\in\mathcal{G}$.\\
Upward closure passes from any member of $\mathcal{C}$ containing $A$. Thus $\mathcal{G}\in\mathscr{P}$ and it bounds $\mathcal{C}$.
\item \hlt{Maximal extension}: Every chain in $\mathscr{P}$ has an upper bound in $\mathscr{P}$, so a maximal $\mathcal{U}\in\mathscr{P}$ exists by Zorn's Lemma.\\
The filter $\mathcal{U}$ is an ultrafilter by Theorem \ref{thm:ultrafilters-are-maximal-filters}, and $\mathcal{F}\subseteq\mathcal{U}$ by the definition of $\mathscr{P}$.
\end{enumerate}
\hlt{Structural Commentary}: Unions preserve filters along chains, so Zorn produces a maximal extension.\\
\hlt{Key Insight}: Chain compatibility makes binary filter closure survive union.
\end{proof}
\end{theorem}

\begin{remark} \label{rem:choice-strength-of-ultrafilter-extension} \hlt{Choice Strength of Ultrafilter Extension}\\
Over ZF, Principle \ref{princ:ultrafilter-lemma} is equivalent to the Boolean Prime Ideal Theorem and is strictly weaker than the full Axiom of Choice.\\
This metamathematical comparison is not used below as proof authority. Theorem \ref{thm:ultrafilter-extension-under-zorn} establishes only the implication from the stronger Zorn principle available earlier. No choice principle enters Theorem \ref{thm:ultrafilters-are-maximal-filters} or the definition of principal ultrafilters.
\end{remark}

\begin{corollary} \label{cor:nonprincipal-ultrafilter-on-infinite-set} \hlt{Nonprincipal Ultrafilters on Infinite Sets}\\
Assume Principle \ref{princ:ultrafilter-lemma}, and let $X$ be infinite. Then $(\exists\mathcal{U}\subseteq\mathcal{P}(X))\,[\mathcal{U}\text{ is a nonprincipal ultrafilter on }X]$.
\begin{proof}
The family $\operatorname{Fr}(X)$ is a filter by Proposition \ref{prop:basic-filter-and-ideal-constructions}(iii). Choose an ultrafilter $\mathcal{U}\supseteq\operatorname{Fr}(X)$ by the hypothesis of the corollary.\\
If $\mathcal{U}$ were principal at $x\in X$, then $\{x\}\in\mathcal{U}$. Since $X\setminus\{x\}\in\operatorname{Fr}(X)\subseteq\mathcal{U}$, closure under intersections would give $\emptyset\in\mathcal{U}$, contradicting properness. Hence $\mathcal{U}$ is nonprincipal.\\
\hlt{Structural Commentary}: Extend the cofinite filter and exclude principality by intersecting complementary singleton data.\\
\hlt{Key Insight}: Cofinite membership prevents concentration at one point.
\end{proof}
\end{corollary}

\begin{definition} \label{def:measurable-cardinal-via-ultrafilter} \hlt{Measurable Cardinal}\\
Let $\kappa$ be an uncountable cardinal. In $\mathrm{ZFC}$, the cardinal $\kappa$ is \hlt{measurable} iff $(\exists\mathcal{U}\subseteq\mathcal{P}(\kappa))\,[\mathcal{U}\text{ is a nonprincipal $\kappa$-complete ultrafilter on }\kappa]$.
\end{definition}

\begin{remark} \label{rem:measurable-cardinal-assumption-boundary} \hlt{Large-Cardinal Assumption Boundary}\\
The assertion that a measurable cardinal exists is an additional large-cardinal hypothesis; it is not assumed elsewhere in this part.\\
The ordinary Ultrafilter Lemma does not supply such a cardinal, as extending a filter to an ultrafilter need not preserve $\kappa$-completeness.\\
Thus the large-cardinal connection enters through the extra completeness requirement, not through ultrafilter maximality alone.
\end{remark}

\begin{assumption} \label{assump:club-and-stationary-setting} \hlt{Club and Stationary Setting}\\
For the club-and-stationary results below, assume Principle \ref{princ:axiom-of-choice-analysis} and fix a regular cardinal $\kappa>\omega$.
\end{assumption}

\begin{definition} \label{def:club-subset-of-regular-cardinal} \hlt{Closed Unbounded Subset}\\
A set $C\subseteq\kappa$ is \hlt{closed in $\kappa$} iff $(\forall\nu<\kappa)\,[\nu\text{ is a nonzero limit ordinal}\land C\cap\nu\text{ is cofinal in }\nu\Rightarrow\nu\in C]$.\\
It is \hlt{unbounded in $\kappa$} iff it is cofinal in $\kappa$, and it is a \hlt{club subset of $\kappa$} iff it is both closed and unbounded.
\end{definition}

\begin{theorem} \label{thm:intersection-of-fewer-than-kappa-clubs} \hlt{Intersection of Fewer Than $\kappa$ Clubs}\\
If $0<\lambda<\kappa$ and $(C_i)_{i<\lambda}$ is a family of club subsets of $\kappa$, then $\bigcap\limits_{i<\lambda}C_i$ is club in $\kappa$.
\begin{proof}
Set $C\defeq\bigcap\limits_{i<\lambda}C_i$. If $\nu<\kappa$ is a nonzero limit ordinal and $C\cap\nu$ is cofinal in $\nu$, then $(\forall i<\lambda)\,C_i\cap\nu$ is cofinal in $\nu$.\\
Hence $(\forall i<\lambda)\,\nu\in C_i$, so $\nu\in C$ and $C$ is closed.\\
If $\theta<\kappa$ and $g:\theta\to\kappa$, then $\sup g(\theta)<\kappa$; otherwise $g$ would be cofinal and $\operatorname{cf}(\kappa)\leq\theta<\kappa$, contradicting regularity.\\
Fix $\beta<\kappa$ and set $\beta_0\defeq\beta+1<\kappa$ by Lemma \ref{lem:infinite-cardinals-are-limit-ordinals}.\\
Recursively, after $\beta_n$ is defined, choose $(\gamma_{n,i})_{i<\lambda}$ with $(\forall i<\lambda)\,\gamma_{n,i}\in C_i\setminus(\beta_n+1)$ by the standing Choice Principle.\\
Set $\beta_{n+1}\defeq\sup\limits_{i<\lambda}(\gamma_{n,i}+1)$. The limit-ordinal property of $\kappa$ gives $(\forall i<\lambda)\,\gamma_{n,i}+1<\kappa$, so the preceding bound gives $\beta_{n+1}<\kappa$.\\
Set $\delta\defeq\sup\limits_{n<\omega}\beta_n$. The same observation with $\theta=\omega$ gives $\delta<\kappa$.\\
Fix $i<\lambda$. Since $(\forall n<\omega)\,\beta_n<\gamma_{n,i}<\beta_{n+1}\leq\delta$, the sequence $(\gamma_{n,i})_{n<\omega}$ is cofinal in the nonzero limit ordinal $\delta$.\\
Closedness of $C_i$ gives $\delta\in C_i$. As $i<\lambda$ was arbitrary, $\delta\in C$; moreover, $\beta<\delta$, so $C$ is unbounded.\\
\hlt{Structural Commentary}: Regularity bounds a countable recursion whose stages visit every club.\\
\hlt{Key Insight}: Each club contributes a cofinal subsequence to one common limit.
\end{proof}
\end{theorem}

\begin{definition} \label{def:club-filter-stationary-and-nonstationary} \hlt{Club Filter, Stationary Sets, and the Nonstationary Ideal}
\begin{enumerate}[label=(\roman*)]
\setlength{\itemsep}{0pt}
\item The \hlt{club filter} on $\kappa$ is $\mathcal{F}_{\mathrm{club}}(\kappa)\defeq\{A\subseteq\kappa:(\exists C\subseteq A)\,C\text{ is club in }\kappa\}$.
\item A set $S\subseteq\kappa$ is \hlt{stationary} iff $(\forall C\subseteq\kappa)\,[C\text{ is club in }\kappa\Rightarrow S\cap C\neq\emptyset]$.
\item The \hlt{nonstationary ideal} is $\operatorname{NS}_{\kappa}\defeq\{A\subseteq\kappa:(\exists C\subseteq\kappa)\,[C\text{ is club in }\kappa\land A\cap C=\emptyset]\}$. Thus $S$ is stationary iff $S\notin\operatorname{NS}_{\kappa}$.
\end{enumerate}
\end{definition}

\begin{proposition} \label{prop:club-filter-completeness-and-stationary-basics} \hlt{Club-Filter Completeness and Stationary-Set Basics}\\
Let $C,S\subseteq\kappa$.
\begin{enumerate}[label=(\roman*)]
\setlength{\itemsep}{0pt}
\item The family $\mathcal{F}_{\mathrm{club}}(\kappa)$ is a $\kappa$-complete filter.
\item The ideal $\operatorname{NS}_{\kappa}$ is dual to $\mathcal{F}_{\mathrm{club}}(\kappa)$.
\item If $C$ is club and $S$ is stationary, then $C\cap S$ is stationary.
\item If $S$ is stationary, then $S$ is unbounded.
\item If $S$ is stationary, then $\{\nu\in S:\nu\text{ is a nonzero limit ordinal}\}$ is stationary.
\end{enumerate}
\begin{proof}
\begin{enumerate}[label=(\roman*)]
\setlength{\itemsep}{0pt}
\item The set $\kappa$ is club in itself, so $\kappa\in\mathcal{F}_{\mathrm{club}}(\kappa)$. No club is empty, so $\emptyset\notin\mathcal{F}_{\mathrm{club}}(\kappa)$, and upward closure follows from the definition.\\
Let $\lambda<\kappa$ and $(A_i)_{i<\lambda}\in\mathcal{F}_{\mathrm{club}}(\kappa)^\lambda$. If $\lambda=0$, then $\bigcap\limits_{i<\lambda}A_i=\kappa$. If $\lambda>0$, choose clubs $C_i\subseteq A_i$ by the standing Choice Principle.\\
The intersection $\bigcap\limits_{i<\lambda}C_i$ is club by Theorem \ref{thm:intersection-of-fewer-than-kappa-clubs}, and it is contained in $\bigcap\limits_{i<\lambda}A_i$. Thus the club filter is $\kappa$-complete.
\item For $A\subseteq\kappa$, $\kappa\setminus A\in\mathcal{F}_{\mathrm{club}}(\kappa)$ iff $(\exists C\subseteq\kappa\setminus A)\,[C\text{ is club}]$, equivalently iff $A\in\operatorname{NS}_{\kappa}$.\\
Hence $\operatorname{NS}_{\kappa}=\mathcal{I}_{\mathcal{F}_{\mathrm{club}}(\kappa)}$, which is an ideal by (i) and Proposition \ref{prop:basic-filter-and-ideal-constructions}(i).
\item Fix a club $D\subseteq\kappa$. The intersection $C\cap D$ is club by Theorem \ref{thm:intersection-of-fewer-than-kappa-clubs} with $\lambda=2$, so $(C\cap S)\cap D\neq\emptyset$.\\
As $D$ was arbitrary, $C\cap S$ is stationary.
\item Suppose $S\subseteq\kappa$ is bounded, and choose $\beta<\kappa$ with $S\subseteq\beta+1$. Set $D_\beta\defeq\{\xi<\kappa:\beta<\xi\}$.\\
Since $\kappa$ is a limit ordinal, $(\forall\eta<\kappa)\,\max\{\eta,\beta\}+1\in D_\beta$, so $D_\beta$ is unbounded.\\
If $D_\beta\cap\nu$ is cofinal in a nonzero limit ordinal $\nu<\kappa$, then $D_\beta\cap\nu\neq\emptyset$, so $\beta<\nu$ and $\nu\in D_\beta$.\\
Hence $D_\beta$ is closed and therefore club. Since $S\cap D_\beta=\emptyset$, the set $S$ is not stationary.
\item Set $L\defeq\{\nu<\kappa:\nu\text{ is a nonzero limit ordinal}\}$. If $L\cap\nu$ is cofinal in a nonzero limit ordinal $\nu<\kappa$, then $\nu\in L$, so $L$ is closed.\\
For $\beta<\kappa$, the sequence $(\beta+n)_{n<\omega}$ is bounded below $\kappa$ by regularity, so $\beta+\omega<\kappa$.\\
Since $\beta+\omega\in L$ lies above $\beta$, the set $L$ is unbounded and therefore club. The set $S\cap L$ is stationary by (iii).
\end{enumerate}
\hlt{Structural Commentary}: Club intersections prove completeness; closed tails and limit ordinals test stationarity.\\
\hlt{Key Insight}: Stationary sets survive every additional club restriction.
\end{proof}
\end{proposition}

\begin{definition} \label{def:diagonal-intersection-and-normal-filter} \hlt{Diagonal Intersection and Normality}
\begin{enumerate}[label=(\roman*)]
\setlength{\itemsep}{0pt}
\item For $(X_\xi)_{\xi<\kappa}\in\mathcal{P}(\kappa)^\kappa$, its \hlt{diagonal intersection} is $\mathop{\triangle}\limits_{\xi<\kappa}X_\xi\defeq\{\beta<\kappa:(\forall\xi<\beta)\,\beta\in X_\xi\}$.
\item A filter $\mathcal{F}$ on $\kappa$ is \hlt{normal} iff $(\forall(X_\xi)_{\xi<\kappa}\in\mathcal{F}^\kappa)\,\mathop{\triangle}\limits_{\xi<\kappa}X_\xi\in\mathcal{F}$.
\end{enumerate}
\end{definition}

\begin{theorem} \label{thm:club-filter-is-normal} \hlt{The Club Filter Is Normal}
\begin{enumerate}[label=(\roman*)]
\setlength{\itemsep}{0pt}
\item If $(C_\xi)_{\xi<\kappa}$ is a sequence of clubs in $\kappa$, then $\mathop{\triangle}\limits_{\xi<\kappa}C_\xi$ is club.
\item The club filter on $\kappa$ is normal.
\end{enumerate}
\begin{proof}
\begin{enumerate}[label=(\roman*)]
\setlength{\itemsep}{0pt}
\item We prove clubness through decreasing refinement, closure, and recursive unboundedness.
\begin{enumerate}[label=(\alph*)]
\setlength{\itemsep}{0pt}
\item \hlt{Decreasing refinement}: Define $(C'_\xi)_{\xi<\kappa}$ by $C'_\xi\defeq\bigcap\limits_{\eta\leq\xi}C_\eta$.\\
Each $C'_\xi$ is club by Theorem \ref{thm:intersection-of-fewer-than-kappa-clubs} with $\lambda=\xi+1$, since $0<\xi+1<\kappa$; moreover, $(C'_\xi)_{\xi<\kappa}$ is decreasing.\\
Fix $\rho<\kappa$. Then $(\forall\xi<\rho)\,\rho\in C'_\xi$ iff $(\forall\xi<\rho)(\forall\eta\leq\xi)\,\rho\in C_\eta$ iff $(\forall\eta<\rho)\,\rho\in C_\eta$.\\
As $\rho<\kappa$ was arbitrary, the diagonal intersections agree. Hence assume $(C_\xi)_{\xi<\kappa}$ is decreasing and set $D\defeq\mathop{\triangle}\limits_{\xi<\kappa}C_\xi$.
\item \hlt{Closedness}: Suppose $\nu<\kappa$ is a nonzero limit ordinal and $D\cap\nu$ is cofinal in $\nu$. Fix $\xi<\nu$.\\
Then $\{\rho\in D\cap\nu:\xi<\rho\}\subseteq C_\xi$ is cofinal in $\nu$, so $\nu\in C_\xi$ by closedness.\\
As $\xi<\nu$ was arbitrary, $\nu\in D$; thus $D$ is closed.
\item \hlt{Unboundedness}: Fix $\eta<\kappa$ and set $\beta_0\defeq\eta+1<\kappa$.\\
Recursively set $\beta_{n+1}\defeq\min(C_{\beta_n}\setminus(\beta_n+1))$; unboundedness of $C_{\beta_n}$ makes the defining set nonempty.\\
Set $\beta\defeq\sup\limits_{n<\omega}\beta_n<\kappa$; otherwise $n\mapsto\beta_n$ would be cofinal in $\kappa$, contrary to $\operatorname{cf}(\kappa)=\kappa>\omega$.\\
The strict increase makes $\beta$ a nonzero limit ordinal. Fix $\xi<\beta$ and choose $n<\omega$ with $\xi<\beta_n$.\\
Decreasingness gives $(\forall m\geq n)\,\beta_{m+1}\in C_{\beta_m}\subseteq C_\xi$. This tail is cofinal in $\beta$, so $\beta\in C_\xi$ by closedness.\\
As $\xi<\beta$ was arbitrary, $\beta\in D$; since $\eta<\beta$, the set $D$ is unbounded and therefore club.
\end{enumerate}
\item Let $(A_\xi)_{\xi<\kappa}\in\mathcal{F}_{\mathrm{club}}(\kappa)^\kappa$, and choose clubs $C_\xi\subseteq A_\xi$ by the standing Choice Principle.\\
The set $\mathop{\triangle}\limits_{\xi<\kappa}C_\xi$ is club by (i), and $\mathop{\triangle}\limits_{\xi<\kappa}C_\xi\subseteq\mathop{\triangle}\limits_{\xi<\kappa}A_\xi$. Hence $\mathop{\triangle}\limits_{\xi<\kappa}A_\xi\in\mathcal{F}_{\mathrm{club}}(\kappa)$.
\end{enumerate}
\hlt{Structural Commentary}: Decreasing refinement separates closure from recursive unboundedness; contained clubs establish normality.\\
\hlt{Key Insight}: Each point satisfies only club conditions indexed below it.
\end{proof}
\end{theorem}

\begin{theorem} \label{thm:fodors-pressing-down-lemma} \hlt{Fodor's Pressing-Down Lemma}\\
Let $S\subseteq\kappa$ be stationary and let $f:S\to\kappa$ be \hlt{regressive}, meaning $(\forall\alpha\in S)\,f(\alpha)<\alpha$. Then $(\exists\gamma<\kappa)\,[\{\alpha\in S:f(\alpha)=\gamma\}\text{ is stationary}]$.
\begin{proof}
For $\gamma<\kappa$, set $S_\gamma\defeq\{\alpha\in S:f(\alpha)=\gamma\}$. Suppose $(\forall\gamma<\kappa)\,[S_\gamma\text{ is nonstationary}]$.\\
Choose $(C_\gamma)_{\gamma<\kappa}$ satisfying $(\forall\gamma<\kappa)\,[C_\gamma\text{ is club}\land C_\gamma\cap S_\gamma=\emptyset]$ by the standing Choice Principle.\\
The set $D\defeq\mathop{\triangle}\limits_{\gamma<\kappa}C_\gamma$ is club by Theorem \ref{thm:club-filter-is-normal}(i). Choose $\alpha\in S\cap D$ and set $\gamma\defeq f(\alpha)$. Then $\gamma<\alpha$ and $\alpha\in D$, so $\alpha\in C_\gamma$.\\
Since $\alpha\in S_\gamma$, this contradicts $C_\gamma\cap S_\gamma=\emptyset$. Hence $(\exists\gamma<\kappa)\,[S_\gamma\text{ is stationary}]$.\\
\hlt{Structural Commentary}: Avoid every fiber by a club, combine the clubs diagonally, and use regressiveness.\\
\hlt{Key Insight}: The regressive value indexes a club already tested at its argument.
\end{proof}
\end{theorem}

\begin{assumption} \label{assump:descriptive-set-theory-setting} \hlt{Descriptive-Set-Theory Setting}\\
The cardinal fixed in Assumption \ref{assump:club-and-stationary-setting} is no longer in force. Work in $\mathrm{ZFC}$, and let $X$ be Baire space, Cantor space, or a finite product of these spaces with its product topology.
\end{assumption}

\begin{samepage}
\begin{definition} \label{def:borel-sigma-algebra-and-pointclasses} \hlt{Borel Sets and the Borel Pointclasses}\\
The \hlt{Borel $\sigma$-algebra} $\mathcal{B}(X)$ is the least family of subsets of $X$ containing all open sets and closed under complements and countable unions. Its members are the \hlt{Borel sets}.\\
Let $\omega_1$ be the least uncountable ordinal. The \hlt{Borel pointclasses} are defined recursively for $0<\alpha<\omega_1$.
\begin{enumerate}[label=(\roman*)]
\setlength{\itemsep}{0pt}
\item $\boldsymbol{\Sigma}^{0}_{1}(X)\defeq\{A\subseteq X:A\text{ is open}\}$ and $\boldsymbol{\Pi}^{0}_{1}(X)\defeq\{X\setminus A:A\in\boldsymbol{\Sigma}^{0}_{1}(X)\}$.
\item Fix $1<\alpha<\omega_1$ after all lower levels have been defined.\\
Set $\boldsymbol{\Sigma}^{0}_{\alpha}(X)\defeq\{\bigcup\limits_{n<\omega}A(n):A\in(\bigcup\limits_{0<\beta<\alpha}\boldsymbol{\Pi}^{0}_{\beta}(X))^{\omega}\}$. Set $\boldsymbol{\Pi}^{0}_{\alpha}(X)\defeq\{X\setminus A:A\in\boldsymbol{\Sigma}^{0}_{\alpha}(X)\}$.
\item $\boldsymbol{\Delta}^{0}_{\alpha}(X)\defeq\boldsymbol{\Sigma}^{0}_{\alpha}(X)\cap\boldsymbol{\Pi}^{0}_{\alpha}(X)$.
\end{enumerate}
Thus $\boldsymbol{\Sigma}^{0}_{1}(X)$, $\boldsymbol{\Pi}^{0}_{1}(X)$, and $\boldsymbol{\Delta}^{0}_{1}(X)$ are the open, closed, and clopen subsets of $X$, respectively.
\end{definition}
\end{samepage}

\begin{theorem} \label{thm:borel-hierarchy-exhausts-borel-sets} \hlt{The Borel Hierarchy Exhausts the Borel Sets}\\
A set $A\subseteq X$ is Borel iff $(\exists\alpha<\omega_1)\,[0<\alpha\land A\in\boldsymbol{\Sigma}^{0}_{\alpha}(X)]$.
\begin{proof}
\begin{enumerate}[label=\roman*.]
\setlength{\itemsep}{0pt}
\item[$(\Rightarrow)$] Set $\mathcal{H}\defeq\bigcup\limits_{0<\alpha<\omega_1}\boldsymbol{\Sigma}^{0}_{\alpha}(X)$. Then $\boldsymbol{\Sigma}^{0}_{1}(X)\subseteq\mathcal{H}$.\\
The relation $(\forall\alpha<\omega_1)\,[0<\alpha\Rightarrow\boldsymbol{\Pi}^{0}_{\alpha}(X)\subseteq\boldsymbol{\Sigma}^{0}_{\alpha+1}(X)]$ holds by constant unions, so $\mathcal{H}$ is complement-closed.\\
Fix $B\in\mathcal{H}^{\omega}$ and choose $\alpha\in(\omega_1\setminus\{0\})^{\omega}$ satisfying $(\forall n<\omega)\,B(n)\in\boldsymbol{\Sigma}^{0}_{\alpha(n)}(X)$ by Choice.\\
One has $(\forall n<\omega)\,B(n)\in\boldsymbol{\Pi}^{0}_{\alpha(n)+1}(X)$ by constant unions and complementation.\\
Set $\gamma\defeq\sup\limits_{n<\omega}(\alpha(n)+1)+1<\omega_1$ by Corollary \ref{cor:at-most-countable-union-at-most-countable}. Since $(\forall n<\omega)\,\alpha(n)+1<\gamma$, one has $\bigcup\limits_{n<\omega}B(n)\in\boldsymbol{\Sigma}^{0}_{\gamma}(X)\subseteq\mathcal{H}$.\\
Thus $\mathcal{H}$ is a $\sigma$-algebra containing $\boldsymbol{\Sigma}^{0}_{1}(X)$, so $\mathcal{B}(X)\subseteq\mathcal{H}$ by minimality.
\item[$(\Leftarrow)$] Set $P(\alpha)\defeq\boldsymbol{\Sigma}^{0}_{\alpha}(X)\cup\boldsymbol{\Pi}^{0}_{\alpha}(X)\subseteq\mathcal{B}(X)$. We prove $(\forall\alpha<\omega_1)\,[0<\alpha\Rightarrow P(\alpha)]$ by strong transfinite induction.\\
The property $P(1)$ holds by the definition of $\mathcal{B}(X)$. Fix $1<\alpha<\omega_1$ and assume $(\forall\beta<\alpha)\,[0<\beta\Rightarrow P(\beta)]$.\\
The inclusion $\boldsymbol{\Sigma}^{0}_{\alpha}(X)\subseteq\mathcal{B}(X)$ holds by the recursive definition, induction hypothesis, and countable-union closure.\\
Thus $\boldsymbol{\Pi}^{0}_{\alpha}(X)\subseteq\mathcal{B}(X)$ by complementation. Hence $P(\alpha)$.
\end{enumerate}
\hlt{Structural Commentary}: Forward closure identifies the generated $\sigma$-algebra; transfinite induction proves the reverse inclusion.\\
\hlt{Key Insight}: Countable Borel operations have ranks below $\omega_1$.
\end{proof}
\end{theorem}

\begin{definition} \label{def:projective-pointclasses} \hlt{The Projective Pointclasses}\\
For $E\subseteq X\times\omega^{\omega}$, set $\pi_X(E)\defeq\{x\in X:(\exists y\in\omega^{\omega})\,(x,y)\in E\}$. Define the projective pointclasses recursively.
\begin{enumerate}[label=(\roman*)]
\setlength{\itemsep}{0pt}
\item $\boldsymbol{\Sigma}^{1}_{1}(X)\defeq\{\pi_X(E):E\in\mathcal{B}(X\times\omega^{\omega})\}$; its members are the \hlt{analytic sets}.
\item For $n\in\N$, $\boldsymbol{\Pi}^{1}_{n}(X)\defeq\{X\setminus A:A\in\boldsymbol{\Sigma}^{1}_{n}(X)\}$.
\item For $n\in\N$, $\boldsymbol{\Sigma}^{1}_{n+1}(X)\defeq\{\pi_X(E):E\in\boldsymbol{\Pi}^{1}_{n}(X\times\omega^{\omega})\}$.
\item For $n\in\N$, $\boldsymbol{\Delta}^{1}_{n}(X)\defeq\boldsymbol{\Sigma}^{1}_{n}(X)\cap\boldsymbol{\Pi}^{1}_{n}(X)$.
\end{enumerate}
The members of $\boldsymbol{\Pi}^{1}_{1}(X)$ are the \hlt{coanalytic sets}. A set $A\subseteq X$ is \hlt{projective} iff $(\exists n\in\N)\,[A\in\boldsymbol{\Sigma}^{1}_{n}(X)\lor A\in\boldsymbol{\Pi}^{1}_{n}(X)]$.
\end{definition}

\begin{proposition} \label{prop:borel-sets-are-delta-one-one} \hlt{Borel Sets Are Analytic and Coanalytic}
\begin{enumerate}[label=(\roman*)]
\setlength{\itemsep}{0pt}
\item $\mathcal{B}(X)\subseteq\boldsymbol{\Delta}^{1}_{1}(X)$.
\item The class $\bigcup\limits_{n\in\N}(\boldsymbol{\Sigma}^{1}_{n}(X)\cup\boldsymbol{\Pi}^{1}_{n}(X))$ is closed under complements.
\end{enumerate}
\begin{proof}
\begin{enumerate}[label=(\roman*)]
\setlength{\itemsep}{0pt}
\item Set $Y\defeq\omega^{\omega}$ and $\mathscr{S}\defeq\{C\subseteq X:C\times Y\in\mathcal{B}(X\times Y)\}$.\\
Then $(\forall C\in\mathscr{S})\,[(X\setminus C)\times Y=(X\times Y)\setminus(C\times Y)\in\mathcal{B}(X\times Y)]$, so $\mathscr{S}$ is complement-closed.\\
One has $(\forall C\in\mathscr{S}^{\omega})\,[(\bigcup\limits_{n<\omega}C(n))\times Y=\bigcup\limits_{n<\omega}(C(n)\times Y)\in\mathcal{B}(X\times Y)]$. Thus $\mathscr{S}$ is countable-union-closed.\\
The inclusion $\boldsymbol{\Sigma}^{0}_{1}(X)\subseteq\mathscr{S}$ holds because products of open sets are open. Thus $\mathscr{S}$ is a $\sigma$-algebra and $\mathcal{B}(X)\subseteq\mathscr{S}$ by minimality.\\
Since $Y\neq\emptyset$, $(\forall C\in\mathcal{B}(X))\,C=\pi_X(C\times Y)\in\boldsymbol{\Sigma}^{1}_{1}(X)$ by the definition of $\boldsymbol{\Sigma}^{1}_{1}(X)$.\\
Consequently, $(\forall A\in\mathcal{B}(X))\,A\in\boldsymbol{\Sigma}^{1}_{1}(X)\cap\boldsymbol{\Pi}^{1}_{1}(X)=\boldsymbol{\Delta}^{1}_{1}(X)$ by the pointclass definitions.
\item Set $\mathscr{P}\defeq(\boldsymbol{\Sigma}^{1}_{n}(X)\cup\boldsymbol{\Pi}^{1}_{n}(X))_{n\in\N}$.\\
Then $(\forall n\in\N)(\forall A\in\mathscr{P}(n))\,X\setminus A\in\mathscr{P}(n)$ by the definition of $\boldsymbol{\Pi}^{1}_{n}(X)$. Thus $\bigcup\limits_{n\in\N}\mathscr{P}(n)$ is complement-closed.
\end{enumerate}
\hlt{Structural Commentary}: Product closure proves the level-one inclusion; pointclass duality proves complement closure.\\
\hlt{Key Insight}: A constant witness realizes every Borel set as a projection.
\end{proof}
\end{proposition}

\begin{remark} \label{rem:descriptive-set-theory-assumption-boundary} \hlt{Determinacy and Large-Cardinal Boundary}\\
The material from Definition \ref{def:borel-sigma-algebra-and-pointclasses} through Proposition \ref{prop:borel-sets-are-delta-one-one} is developed in $\mathrm{ZFC}$ alone, without measurable-cardinal or determinacy hypotheses. No projective-regularity theorem is asserted; any later theorem of that kind must state its additional hypothesis explicitly.
\end{remark}
